\chapter{Values in the code}\label{ch:code}

While the debugger is paused, the code itself shows what the script computed: in the tooltips, at the
ends of the lines, and as a preview in Rhino of whatever is under the mouse. You read the values where
you read the code, without opening the Variables tab or adding watches name by name.

\section{Tooltips}\label{sec:tooltips}

Hover a name in the code. The tooltip has three parts:

\begin{description}
  \item[Its value, as a line of code,] syntax-colored and compact: \code{count = 5},
        \code{pts = [[0,0,0],[5,0,0],… +2]}, \code{p = (1,2,3)}, \code{node = Node(x=3,label='n1')},
        \code{def scale(v,k=2)}. Long collections show their first items and how many more:
        \code{many = [0,1,2,… +247]}.
  \item[Facts, in gray,] for collections and Rhino geometry: \emph{250 items};
        \emph{closed · degree 3 · length 31.4 · 12 control points}; \emph{solid · 6 faces · volume 1000};
        \emph{closed · 482 vertices · 960 faces}. After a step, \emph{was 4} when the value changed.
  \item[The editor's own text,] below a line: the type and its documentation.
\end{description}

Where a name has no value (not paused, not in scope) only the editor's own tooltip shows.

\paragraph{Scope.} Names are read as Python sees them at the hovered line (\see{sec:frame}). A global
\code{crv} and a function's local \code{crv} show their own values. A local of a function that isn't
running shows nothing.

\section{Links in the tooltip}\label{sec:tooltiplinks}

Under the value, the tooltip offers what can be done with it:

\begin{xltabular}{\linewidth}{@{}P{0.28\linewidth}L@{}}
\toprule
\ui{Watch in Rhino} & Adds it to \vn{Watch} (anything but plain numbers and text). \\
\ui{Edit in Rhino}  & Turns it into Rhino objects you can edit (geometry that can be written back, \see{sec:editinrhino}). \\
\ui{Bake}           & Copies it into the document (\see{sec:bake}). \\
\ui{Tree}           & Opens the value's members in a small window (below). \\
\ui{Watch rg.Point3d(…)} & Over a class called on the line, such as \code{Point3d} in
  \code{pts.append(rg.Point3d(ri * math.cos(th), 0, h))}: the whole call, the value it makes, onto
  \vn{Watch}, computed again at every pause. Offered only for classes: a watch of \code{pts.append(…)}
  would append at every pause. For any other expression, select it and right-click \ui{Watch in Rhino}. \\
\bottomrule
\end{xltabular}

\paragraph{The Tree window} opens at the mouse with the value's members, expandable like the Variables
tab. It stays open while you step and reads the value again at every pause, keeping the rows you opened.
It is grayed while the script runs, and closes when the run ends or with \key{Esc}.

\section{Hover preview}\label{sec:hoverpreview}

The geometry of the name under the mouse, such as \code{crv}, \code{crv.PointAtStart} or \code{curves[0]}, is
drawn in yellow in Rhino until the mouse moves on. It stays while the mouse is on the tooltip, so you can
read the values and click its links, and while the code's context menu is open, so \ui{Watch in Rhino}
can take it from there. Setting: \setting{Hover preview}; command \code{miniVivisection HoverPreview=Off}.

\section{Values at the ends of lines}\label{sec:inlinevalues}

Each line of the paused function, from its start to the paused line, shows the values of the variables
written on it, in gray after the line: \code{total = 6, k = 3}, as in VS~Code and PyCharm.

After a step, values that changed are orange, and the tooltip says what they were before
(\emph{was 5}). The comparison is within the same function, and per thread: stepping into another
function starts afresh. Settings: \setting{Values at the ends of lines},
\setting{What changed since the last pause}.

\section{The editor's Watch tab evaluates expressions}\label{sec:watchtab}

The Script Editor's own Watch tab only resolves plain variable names and shows anything else as an
error. With miniVivisection it evaluates any Python expression in the paused frame: \code{1+2} shows
\code{3}, \code{crv.PointAtStart} its point, \code{[c.PointAtEnd for c in curves]} a list to expand. A Python
error shows as the row's value (\code{NameError: name 'x' is not defined}). This tab isn't drawn: drawing is
\vn{Watch}'s job, and the two lists are separate. Setting: \setting{Expressions in the editor's Watch tab};
command \code{miniVivisection WatchExpressions=Off}. C\# keeps plain names, and an outer frame chosen in the
Call Stack tab keeps the editor's own lookup.

\section{Live Variables and Watch tabs}

After a value is written back from Rhino, after a \vn{Python} command, or after Claude changed
something, the editor's own Variables and Watch tabs show the new values right away; expanded rows stay
expanded. A moved curve's end point coordinates update while still paused.

\section{Context menus while paused}\label{sec:contextmenus}

Right-click a name in the code (the items act on the selection, or else on the name under the cursor),
or a row in the Variables or Watch tab (the items act on that row's path):

\begin{itemize}
  \item \ui{Watch in Rhino}, and \ui{Watch in Rhino as ▸} with a chosen interpretation (\see{sec:interpretas});
  \item \ui{Bake to Rhino};
  \item \ui{Edit in Rhino}.
\end{itemize}

The menus show only what can be used right now. While paused, most of the code view's own items are
disabled, since the script is read-only then.

\begin{lesson}{lesson 02}
\lessonfile{02_values_in_the_code.py}: tooltips, values at line ends, \emph{was}, the Tree; stepping with \key{F10}.
\end{lesson}
